\completa{6}{Arbués}{3a: de una fila a los extremos por sexo}
\senalar{Recorrer los cuatro cuadros de izquierda a derecha; después, relacionar cada fragmento de código con su cuadro.}
\decir{Esta pregunta nace de dos campos: SEXO permite formar categorías y CANTIDAD permite comparar sus conteos. El punto tres a pide agrupar una categoría y obtener los valores mayores y menores de un campo numérico. Por eso preguntamos: ¿cuál es el mayor y el menor conteo por registro para cada sexo, y dónde se observa un máximo? Buscamos extremos de las filas agregadas.}
\decir{El primer cuadro dice «Entrada: sexo y cantidad». Partimos de las noventa y tres mil ochocientas treinta y una filas; conservamos también territorio, año y mes para localizar extremos. La primera flecha conduce a «uno, Map: sexo, par c y c». Aquí c significa cantidad. El Map emite una clave, el sexo, y una tupla, es decir, un conjunto ordenado de valores. Esa tupla repite la cantidad porque inicia dos cálculos: máximo y mínimo.}
\decir{La siguiente flecha llega a «dos, Reduce: máximo y mínimo por sexo». ReduceByKey combina solamente valores de la misma clave. En el código, h representa una fila y la flecha escrita como igual y mayor que indica qué transformación se aplica a esa fila. En la reducción, a y b son dos tuplas que se están combinando; a punto guion bajo uno y b punto guion bajo uno acceden a sus primeras posiciones, y guion bajo dos, a las segundas. Math punto max conserva el mayor de los primeros valores y math punto min conserva el menor de los segundos. Al repetir la combinación obtenemos un par de extremos por sexo. No sumamos cantidades: elegimos los extremos.}
\decir{La última flecha lleva a «Salida: extremos y ubicación». Las otras reducciones del script conservan una tupla más amplia con cantidad, lugar y fecha. La tabla muestra F: máximo dieciocho y mínimo uno; M: máximo catorce y mínimo uno; y NO ESPECIFICA: máximo tres y mínimo uno. Un máximo F está en Villa María del Triunfo, noviembre de 2019; uno M está en Santiago, Cusco, octubre de 2019. Si hay empate, se conserva un registro extremo; no enumeramos todos los empatados. Ya sabemos los límites: ahora explicaremos qué ocurre entre ellos.}

\completa{7}{Arbués}{3b: cómo se construyen los estadísticos}
\senalar{Leer los cuatro cuadros y las tres filas Map/Reduce; señalar cada símbolo de las fórmulas.}
\decir{Conocer un máximo no nos dice cuál es el valor típico. Del punto tres b y de la columna CANTIDAD nace esta pregunta: ¿cuál es el promedio, la mediana, la desviación, el máximo y el mínimo por registro? La entrada sigue siendo toda la tabla, sin filtrar personas ni departamentos.}
\decir{El primer cuadro dice «Entrada: columna CANTIDAD». La flecha llega a «uno, Map: par c y uno, c al cuadrado y c». Es un resumen de varias ramas del código. En la primera, cada fila emite su cantidad junto con un uno; ese uno contará la fila. En la segunda emitimos la cantidad al cuadrado. En las ramas de extremos emitimos directamente la cantidad.}
\decir{La segunda flecha conduce a «dos, Reduce: sumas, conteo y extremos». La tabla inferior precisa qué combinamos. El par c y uno se reduce sumando componente por componente: S es la suma de cantidades y n es la suma de unos, o número de registros. El símbolo sigma, parecido a una E mayúscula, indica sumatoria. Los cuadrados se suman para obtener Q. Finalmente, las cantidades se combinan con máximo y con mínimo en reducciones separadas. En el fragmento, h representa una fila; toDouble convierte la cantidad a número con decimales y uno L indica un entero de tipo largo.}
\decir{La última flecha llega a «Salida: estadísticos». La media, escrita equis barra, es S dividido entre n. La desviación poblacional, escrita sigma, es la raíz cuadrada de Q dividido entre n menos la media al cuadrado. Resume la variación de los conteos alrededor de la media. La mediana requiere otro procedimiento: ordenar las cantidades, asignarles índices y tomar el centro. Con un número impar de filas, ese centro corresponde a una sola posición. No llamamos reducción al ordenamiento.}
\decir{Leamos la salida: noventa y tres mil ochocientas treinta y una filas; suma de ciento veintiún mil cuatrocientas veintisiete; promedio de uno coma dos nueve cuatro uno; mediana uno; desviación poblacional cero coma ocho nueve cero seis; mínimo uno y máximo dieciocho. La mediana uno y la media superior a uno muestran que los conteos altos elevan el promedio. Media y promedio son sinónimos; mediana es otro estadístico. Con esta base pasaremos de cantidades a proporciones.}

\completa{8}{Arbués}{3c.1: menores, denominador territorial y control nacional}
\senalar{Leer A y su flecha; B y su flecha; C y su flecha. Después señalar el cálculo que une A con B y la tabla.}
\decir{El punto tres c solicita resultados decimales. Aquí usamos edad, departamento y cantidad para preguntar: ¿qué porcentaje del total de cada departamento corresponde a menores de dieciocho años y cuál es el control nacional? El porcentaje permite comparar composición, porque un departamento pequeño puede tener menos menores en número absoluto y una proporción mayor. En el diagrama, d significa departamento y c, cantidad. Menores incluye cero a cinco, seis a once y doce a diecisiete años.}
\decir{En la primera fila leemos «uno, Map, A: total; par d y c; todas las filas». Cada fila emite su departamento como clave y su cantidad como valor. La flecha conduce a «dos, Reduce, A: total; reduceByKey, suma; par d y T sub d, total por departamento». El guion bajo más guion bajo del código significa sumar dos valores que se están combinando. T sub d es el total de personas contabilizadas para ese departamento.}
\decir{La segunda fila dice «uno, Map, B: menores; par d y c; filtrar menores antes del Map». Primero comprobamos el rango de edad; solo las filas de menores emiten ese par. Su flecha llega a «dos, Reduce, B: menores; reduceByKey, suma; par d y M sub d, menores por departamento». M sub d es el numerador. A y B parten del mismo RDD con filtros diferentes: la segunda fila no recibe como entrada el resultado de la primera.}
\decir{La tercera fila dice «uno, Proyección, C: control; values de A y B; tomar los totales parciales». Values extrae los valores y deja de lado las claves departamentales. Su flecha conduce a «dos, Reduce, C: control; reduce, suma; total y menores nacionales». Se reduce la colección de totales de A y, por separado, la de menores de B. Este control sí utiliza resultados previos. Se llama proyección porque el script usa values; no mostramos un Map explícito que no existe en esa línea.}
\decir{Después de reducir, join une A y B por departamento. Para cada clave obtenemos total y menores; entonces un map calcula p sub d igual a cien por M sub d dividido entre T sub d. P es porcentaje, M es numerador y T es denominador. La multiplicación por cien lo expresa por cada cien personas contabilizadas; el cálculo usa decimales y se redondea a dos posiciones. El fragmento inferior corresponde al Map y Reduce de A.}
\decir{En la salida, Madre de Dios tiene cuatrocientos setenta y tres menores de seiscientas cincuenta y cinco personas: setenta y dos coma veintiuno por ciento. Ucayali tiene mil quinientos cuatro de dos mil ciento veintiséis: setenta coma setenta y cuatro por ciento. Lima Metropolitana registra diecinueve mil setecientos cincuenta y uno de treinta y nueve mil cuatrocientos noventa: cincuenta coma cero dos por ciento. El control nacional divide setenta mil ciento setenta y ocho entre ciento veintiún mil cuatrocientas veintisiete y obtiene cincuenta y siete coma setenta y nueve por ciento. El denominador incluye edades no especificadas; no es la población del departamento. Aplicaremos ahora el mismo principio a nacionalidad y año.}

\completa{9}{Arbués}{3c.2: extranjeros por año y apoyo por nacionalidad}
\senalar{Seguir horizontalmente los tres pares; distinguir la clave año de A/B y la clave nacionalidad de C.}
\decir{La nacionalidad y el año permiten plantear otra consulta decimal: ¿qué porcentaje anual corresponde a nacionalidades extranjeras y cuáles aportan más al acumulado? Queremos separar el tamaño del conteo de su peso relativo. En la entrada, a significa año; n, nacionalidad; y c, cantidad. Excluimos PERÚ y NO ESPECIFICA para definir el subconjunto extranjero.}
\decir{La fila A dice «uno, Map, A: total; par a y c; todas las filas». Emitimos año y cantidad. La flecha lleva a «dos, Reduce, A: total; reduceByKey, suma; par a y T sub a, total anual». Así obtenemos el denominador de cada año. En la fila B leemos «uno, Map, B: extranjeros; par a y c; solo extranjeros». Antes de emitir filtramos la nacionalidad. Su flecha conduce a «dos, Reduce, B: extranjeros; reduceByKey, suma; par a y E sub a, extranjeros por año». E sub a es el numerador anual.}
\decir{En la fila C leemos «uno, Map, C: apoyo; par n y c; solo extranjeros». Ahora cambia la clave: agrupamos por nacionalidad, no por año. Su flecha llega a «dos, Reduce, C: apoyo; reduceByKey, suma; total por nacionalidad». Esa rama permite describir qué nacionalidades reúnen más personas en todo el periodo. No es una tercera reducción consecutiva del porcentaje anual. El fragmento de código visible muestra la construcción del total anual de A.}
\decir{Una vez reducidos A y B, los unimos por año y calculamos p sub a igual a cien por E sub a dividido entre T sub a. La salida de 2019 muestra cuatrocientas treinta y cuatro personas extranjeras de veintiséis mil trescientas treinta y una: uno coma sesenta y cinco por ciento. En 2023 son quinientas veinticinco de dieciocho mil ciento noventa y dos: dos coma ochenta y nueve por ciento. En 2024 son quinientas veintiocho de diecinueve mil cuatrocientas veinte: dos coma setenta y dos por ciento.}
\decir{El apoyo C registra dos mil ciento cincuenta y una personas de nacionalidad venezolana y ciento veinticuatro de nacionalidad colombiana. Finalmente, dos coma setenta y dos menos dos coma ochenta y nueve da menos cero coma diecisiete puntos porcentuales. Aunque el conteo extranjero crece de quinientos veinticinco a quinientos veintiocho, su participación baja porque también cambia el total anual. Puntos porcentuales es la resta entre porcentajes; no una tasa de crecimiento del conteo. A continuación trasladaremos esta lógica a los meses.}

\completa{10}{Arbués}{3c.3: participación mensual, referencia e índice}
\senalar{Leer A, B y C con sus flechas; después las dos fórmulas. Aclarar que la agrupación por mes reúne los seis años.}
\decir{A partir de MES y CANTIDAD preguntamos: ¿qué meses superan una distribución uniforme y cómo se concentra el total por trimestre? Esta consulta responde al punto tres c con porcentajes e índices. Reunimos los seis años: enero significa todos los eneros del periodo. La referencia uniforme divide cien entre doce y obtiene ocho coma treinta y tres por ciento aproximadamente para cada mes. Es una referencia descriptiva, no una prueba de estacionalidad.}
\decir{La fila A dice «uno, Map, A: mes; par m y c; clave mes». M representa el número del mes y c, cantidad. La flecha llega a «dos, Reduce, A: mes; reduceByKey, suma; par m y T sub m, total mensual». Combinamos todas las cantidades del mismo mes. El fragmento de Scala muestra precisamente esa emisión y esa suma por clave.}
\decir{En B leemos «uno, Map, B: nacional; par TOTAL y c; la misma clave para todas las filas». TOTAL es una etiqueta constante. Como todas las filas comparten esa clave, la flecha lleva a «dos, Reduce, B: nacional; reduceByKey, suma; T igual a ciento veintiún mil cuatrocientas veintisiete». Obtenemos un único denominador para el periodo completo.}
\decir{La fila C dice «uno, Map, C: apoyo; par t y c; calcular trimestre desde el mes». T minúscula representa trimestre: enero a marzo es el primero, abril a junio el segundo, julio a septiembre el tercero y octubre a diciembre el cuarto. Su flecha llega a «dos, Reduce, C: apoyo; reduceByKey, suma; par t y S sub t, total trimestral». Esta agrupación reúne el mismo trimestre de los seis años; no produce los veinticuatro trimestres cronológicos de otra consulta.}
\decir{Después calculamos participación p sub m igual a cien por T sub m dividido entre T. El índice I sub m es doce por T sub m dividido entre T, equivalente a dividir la participación observada entre la uniforme. Un índice uno coincide con la referencia; mayor que uno la supera y menor que uno queda por debajo. No multiplicamos el porcentaje por doce sin ajustar su escala.}
\decir{La salida muestra diciembre: once mil quinientas veintidós personas, nueve coma cuarenta y nueve por ciento e índice uno coma catorce. Febrero: once mil ciento dieciocho, nueve coma dieciséis por ciento e índice uno coma diez. Abril: ocho mil ochocientas cinco, siete coma veinticinco por ciento e índice cero coma ochenta y siete. El cuarto trimestre acumula treinta y dos mil seiscientas cincuenta, veintiséis coma ochenta y nueve por ciento. Diciembre queda aproximadamente catorce por ciento por encima de la referencia uniforme; eso no significa que su participación sea catorce por ciento. Tampoco ajustamos por días del mes ni demostramos causas. Concluido el bloque RDD, doy paso a Sergio, quien continuará con las consultas sobre tablas con esquema.}
